

Inserir um percentual na frente para comentar e não aparecer na compilação do documento final. 

No arquivo pre\_projeto.tex --- "\ input{exemplos}"

Texto considerando a revisão da literatura pertinente, dividido em seções e subseções.

Este é um exemplo de como usar figuras. Referência cruzada: Figura~\ref{fig:exemplo}


\begin{figure}[!htbp]
	\centering
	\caption{Exemplo de figura}
	%scale redimensiona a figura.
	%1.5 = 150% do tamanho original
	%1 = 100% do tamanho original
	%0.20 = 20% do tamanho original
	\includegraphics[scale=1]{imagens/IFSP-BRA.png}
	\legend{Fonte: Disponível em: http://bra.ifsp.edu.br. Acesso em 27 mar. 2020}.
	\label{fig:exemplo}
\end{figure}



Este é um exemplo de como usar tabelas. Referência cruzada: Tabela~\ref{tab:exemplo}


\begin{table}[!htbp]
\centering
\caption{Exemplo de tabela de 2 colunas}
	\begin{tabular}{ c | c }
		\hline
		\textbf{Coluna 1} & \textbf{Coluna 2} \\ \hline
		Dado 1a           & Dado 2a           \\ \hline
		Dado 1b           & Dado 2b           \\ \hline
		Dado 1c           & Dado 2c           \\ \hline
		Dado 1d           & Dado 2d           \\ \hline
	\end{tabular}
	\\ \vspace{0.2cm}
	Fonte: Autoria própria
	\label{tab:exemplo}
\end{table}



Este é um exemplo de como usar quadros. Referência cruzada: Quadro~\ref{tab:exemploquad}


\begin{quadro}[!htbp]
\centering
\caption{Exemplo de Quadro de 3 colunas}
	\begin{tabular}{ | m{10em} | m{4cm}| m{4cm} | }
		\hline
		\textbf{Coluna 1} & \textbf{Coluna 2} & \textbf{Coluna 3} \\ \hline
		Dado 1a           & Dado 2a & \\ \hline
		Dado 1b           & Dado 2b & \\ \hline
		Dado 1c           & Dado 2c & \\ \hline
		Dado 1d           & Dado 2d & \\ \hline
	\end{tabular}
	\\ \vspace{0.2cm}
	Fonte: Autoria própria
	\label{tab:exemploquad}
\end{quadro}


Este é um exemplo de como usar quadros. Referência cruzada: Quadro~\ref{tab:exemplo2}


\begin{quadro}[!htbp]
\centering
\caption{Exemplo de Quadro de 2 colunas}
	\begin{tabular}{ | m{10em} | m{4cm}| }
		\hline
		\textbf{Coluna 1} & \textbf{Coluna 2}  \\ \hline
		Dado 1a           & Dado 2a  \\ \hline
		Dado 1b           & Dado 2b  \\ \hline
		Dado 1c           & Dado 2c  \\ \hline
		Dado 1d           & Dado 2d  \\ \hline
	\end{tabular}
	\\ \vspace{0.2cm}
	Fonte: Autoria própria
	\label{tab:exemplo2}
\end{quadro}


Este é um exemplo de como usar equações. Referência cruzada: Equação~\ref{eq:exemplo}

\begin{equation}
\sum_{i=1}^{n} i = \frac{n(n+1)}{2}
\label{eq:exemplo}
\end{equation}

\clearpage


Este é um exemplo de como inserir texto sem formatação (ambiente verbatim):

\begin{verbatim}
	Texto sem formatação, como espaçamento igual.
\end{verbatim}


Exemplo de lista de itens:

\begin{itemize}
	\item \textbf{Item 1:} texto...;
	\item \textbf{Item 2:} texto...;
    \begin{itemize}
            \item \textbf{Subitem:} texto...;
            \item \textbf{Subitem:} texto...;
            \item \textbf{Subitem:} texto...;
        \end{itemize}
	\item \textbf{Item 3:} texto...;
	\item \textbf{Item n:} texto....
\end{itemize}


Exemplo de lista numerada:

\begin{enumerate}
	\item \textbf{Item:} texto...;
	\item \textbf{Item:} texto...;
    \begin{enumerate}
        \item \textbf{Subitem:} texto...;
        \item \textbf{Subitem:} texto...;
        \item \textbf{Subitem:} texto...;
    \end{enumerate}
	\item \textbf{Item:} texto...;
	\item \textbf{Item:} texto....
\end{enumerate}

Tipos de referência a serem utilizados no arquivo referencias.bib:
Exemplos de estruturas.\\

@article{<citation key>,
    author        = {},
    title         = {},
    journaltitle  = {},
    year          = {}
}

@online{<citation key>,
    author        = {},
    title         = {},
    year          = {},
    url           = {},
    urlaccessdate = {}
}

@book{<citation key>,
    author        = {},
    title         = {},
    year          = {}
}

@misc{<citation key>,
    author        = {},
    title         = {},
    year          = {},
    url           = {},
    urlaccessdate = {}
}

Exemplos de Notas de rodapé de leis: \\

... instituído pela Lei Complementar Nº 128\footnote{BRASIL. LEI COMPLEMENTAR Nº 128, de 19 de dezembro de 2008. Altera a Lei Complementar Nº 123, de 14 de dezembro de 2006, altera as Leis nos 8.212, de 24 de julho de 1991, 8.213, de 24 de julho de 1991, 10.406, de 10 de janeiro de 2002 – Código Civil, 8.029, de 12 de abril de 1990, e dá outras providências. \textbf{Diário Oficial [da] República Federativa do Brasil}, Brasília, DF, 19 dez. 2008. Disponível em: http://www.planalto.gov.br/ccivil\_03/leis/lcp/lcp128.htm. Acesso em: 24 nov 2021.}
\\
\\

Exemplos de comandos para texto e referências:

\begin{itemize}
	\item Para iniciar um novo parágrafo, basta deixar uma linha em branco no código fonte;
	\item Não force o compilador a pular mais de uma linha, pois terá influência negativa na composição do documento;
	\item Sempre deixe o \LaTeX\ realizar a formatação de parágrafos e posicionamento de elementos;
	\item Utilização de aspas simples (abertura \verb|`|, fechamento \verb|'|): `Texto entre aspas simples';
	\item Utilização de aspas duplas (abertura \verb|``|, fechamento \verb|''|): ``Texto entre aspas duplas'';
	\item Negrito (comando \verb|\textbf|): \textbf{texto em negrito};
	\item Itálico (comando \verb|\textit|): \textit{texto em itálico};
	\item Sublinhado (comando \verb|\underline|): \underline{texto sublinhado};
	\item Negrito e itálico (usar comandos juntos): \textbf{\textit{texto em negrito e itálico}};
	\item Alterar cor do texto (comando \verb|\textcolor{cor}{texto}|):
	\begin{itemize}
		\item Exemplo \verb|\textcolor{red}{texto}|: \textcolor{red}{texto vermelho};
		\item Exemplo \verb|\textcolor[RGB]{255, 102, 0}|: \textcolor[RGB]{255, 102, 0}{texto laranja};
		\item Exemplo \verb|\textcolor[HTML]{006AD7}|: \textcolor[HTML]{006AD7}{texto azul};
	\end{itemize}
	\item Ambiente matemático inline (comando \verb|$ expressão $|): $s = x^2-2x +1$;
	\item Referência normal (comando \verb|\cite|):
	\begin{itemize}
		\item \cite{Agaisse1995};
		\item \cite{Abedi2014};
		\item \cite{Baum2016};
	\end{itemize}
	\item Referência normal com mais de uma obra (comando \verb|\cite|):
	\begin{itemize}
		\item \cite{Abedi2014, Agaisse1995};
		\item \cite{AgapitoTenfen2014, Baum2016, Nelson2014};
	\end{itemize}
	\item Referência nome e ano (comando \verb|\citeauthorandyear|):
	\begin{itemize}
		\item \citeauthorandyear{bervian2007a};
		\item \citeonline{bervian2007a};
		\item \citeauthorandyear{documento2018};
		\item \citeauthorandyear{Abedi2014};
	\end{itemize}
	\item Referência apud ({mais antigo} {mais novo}):\\
	\verb|\apud[pagina]{autor-indireto}{autor-direto}|):
	\begin{itemize}
		\item \apud{Agaisse1995}{Abedi2014};
		\item \apudonline{Agaisse1995}{Abedi2014};
	\end{itemize}
	\item Referência apud [pagina]({mais antigo} {mais novo}):\\
	\verb|\apud[pagina]{autor-indireto}{autor-direto}|):
	\begin{itemize}
		\item \apud[p. 81]{Agaisse1995}{Abedi2014};
		\item \apudonline[p. 81]{Agaisse1995}{Abedi2014};
	\end{itemize}
	
	\item Referência do mesmo autor, mesmo ano com obras distintas 
	\verb|\citeauthorandyear|):
	\begin{itemize}
		\item \citeauthorandyear{Agaisse1996a};
		\item \citeauthorandyear{Agaisse1996b};
		\item \cite{Agaisse1996a};
		\item \cite{Agaisse1996b};
	\end{itemize}
	
	
\end{itemize}


Exemplo 1 de citação direta:

\begin{citacao}
	Os 20 aminoácidos usualmente encontrados como resíduos em proteínas contém um grupo $\alpha$-carboxil, um grupo $\alpha$-amino e um grupo R distinto substituído no átomo de carbono $\alpha$. O átomo de carbono $\alpha$ de todos os aminoácidos, com exceção da glicina, é assimétrico e, portanto, os aminoácidos podem existir em pelo menos duas formas estereoisoméricas. Somente os estereoisômeros L, com uma configuração relacionada à configuração absoluta da molécula de referência L-gliceraldeído, são encontrados em proteínas \cite[p. 81]{Nelson2014}.
\end{citacao}

Exemplo 2 de citação direta:

\begin{citacao}
	\textit{These various insecticidal proteins are synthesized during the stationary phase and accumulate in the mother cell as a crystal inclusion which can account for up to 25\% of the dry weight of the sporulated cells. The amount of crystal protein produced by a B. thuringiensis culture in laboratory conditions (about 0.5 mg of protein per ml) and the size of the crystals (24) indicate that each cell has to synthesize $10^6$ to $2 \times 10^6$ $\delta$-endotoxin molecules during the stationary phase to form a crystal} \cite[p. 1]{Agaisse1995}.
\end{citacao}

Exemplo de nota de rodapé\footnote{Essa é uma nota de rodapé!}.
